% Generated draft: requires mathematical and source review.
\documentclass[11pt]{article}
\usepackage[margin=25mm]{geometry}
\usepackage{fontspec}
\setmainfont{Latin Modern Roman}
\usepackage{amsmath,amssymb,mathtools,unicode-math}
\setmathfont{Latin Modern Math}
\usepackage{enumitem,needspace,xurl,hyperref}
\setlength{\parindent}{0pt}
\setlength{\parskip}{5pt}
\setlength{\emergencystretch}{3em}
\newcommand{\N}{\mathbb{N}}
\newcommand{\Z}{\mathbb{Z}}
\newcommand{\Q}{\mathbb{Q}}
\newcommand{\R}{\mathbb{R}}
\newcommand{\C}{\mathbb{C}}
\newcommand{\F}{\mathbb{F}}
\newcommand{\A}{\mathbb{A}}
\newcommand{\PP}{\mathbb{P}}
\newcommand{\E}{\mathbb{E}}
\newcommand{\eps}{\varepsilon}
\newcommand{\dd}{\,\mathrm{d}}
\begin{document}
\section*{UnsolvedMath ID: 30003089}
\textbf{Source problem code:} catalogue entry
\subsection{Definitions and mathematical statement}
For a measurable space \(\mathcal X\) and a positive-definite kernel \(k\) with reproducing kernel Hilbert space \(\mathcal H_k\), suppose the kernel mean embeddings \(\mu_P=\int k(x,\cdot)\,\dd P(x)\) and \(\mu_Q\) exist. The maximum mean discrepancy is
\[
\operatorname{MMD}_k(P,Q)=\|\mu_P-\mu_Q\|_{\mathcal H_k}.
\]
This definition does not specify the kernel, the sampling experiment, or the distribution classes intended by the source. In particular, the record does not state conditions ensuring that this discrepancy distinguishes all distributions under consideration.

To describe the requested comparison abstractly, let \(\mathcal E_n=\{\mathbb P_{P,n}:P\in\mathcal P\}\) be a statistical experiment indexed by distributions \(P\), with sample-size index \(n\in\N\). For a specified null class \(\mathcal P_0\subseteq\mathcal P\), separation functional \(\Delta(P,\mathcal P_0)\), and alternative class \(\mathcal P_1\subseteq\mathcal P\), set
\[
\mathcal A(r)=\{P\in\mathcal P_1:\Delta(P,\mathcal P_0)\ge r\}.
\]
A randomized test is a measurable function \(\varphi_n\) into \([0,1]\), representing rejection probability. For fixed \(\alpha,\beta\in(0,1)\) with \(\alpha+\beta<1\), define the minimax separation threshold by
\[
r_n^*(\alpha,\beta)=\inf\left\{r>0:\exists\varphi_n\ \text{such that}\quad
\sup_{P\in\mathcal P_0}\mathbb E_{P,n}\varphi_n\le\alpha,
\quad
\sup_{P\in\mathcal A(r)}\mathbb E_{P,n}(1-\varphi_n)\le\beta
\right\}.
\]
The corresponding threshold for a specified standard MMD test is defined by restricting the test in this expression to that procedure. This is a framework for making the question precise, not a claim that the source specifies these choices.

\textbf{Source question.} Determine to what extent the reported gap between an \(n^{-1/2}\) detection boundary for the standard MMD goodness-of-fit test and an \(n^{-1}\) boundary described as typical of parametric tests reflects intrinsic difficulty of nonparametric testing, and to what extent it reflects suboptimality of the standard MMD statistic. For a fixed nonparametric experiment and separation convention, this requires comparing the standard procedure's threshold with the minimax threshold over all tests; comparison with a parametric experiment additionally requires explicitly specified, comparable alternatives and separation measures.

\textbf{Unresolved qualification.} The record supplies neither the experiments and distribution classes nor the separation measure, kernel, standard statistic, calibration, or error criterion underlying the quoted rates. It does not say whether those rates measure MMD, squared MMD, or another discrepancy. Squaring a discrepancy changes its rate exponent, so the two quoted powers cannot by themselves establish a statistical efficiency gap. Consequently, the supplied question does not determine a unique quantitative theorem or conjecture; the abstract framework above leaves its missing specifications unresolved.
\subsection{Short English statement}
The question asks how much of the reported \(n^{-1/2}\) versus \(n^{-1}\) detection gap comes from the difficulty of nonparametric testing and how much comes from the standard MMD test. A precise comparison requires specified models, alternatives, error bounds, and a common separation convention. These details are absent from the supplied record, so the quoted rates are retained as source claims requiring clarification.
\subsection{Sources}
\begin{itemize}
\item Local UnsolvedMath registry, permanent ID \(30003089\), title ``Detection Gaps Between Parametric and Kernel Tests,'' exact source problem code \texttt{catalogue entry}. The record attributes the question to a 2016 Oberwolfach workshop report; no individual proposer is supplied. Its original formulation asks how much of the gap between the relevant statistical procedures is attributable to the fundamental difference between parametric and nonparametric testing problems.
\item Supplied report citation: \emph{Computationally and Efficient Inference for Complex Large-scale Data} (2016), Oberwolfach Reports, DOI 10.4171/owr/2016/16, pp.\ 741--796. Supplied literature-source label: ``OWR report 16/2016.'' \url{https://doi.org/10.4171/owr/2016/16}. The supplied background describes Oberwolfach Reports as records of mathematical workshops and their open-problem discussions; no passage-level page locator is supplied here.
\item Supplied literature-source label: ``Balasubramanian et al.'' \url{https://www.jmlr.org/papers/volume22/17-570/17-570.pdf}. The record supplies no full bibliographic citation or theorem locator for this reference.
\item Literature assessment dated exactly \texttt{2026-08-21}: the supplied curation labels the problem ``partially solved'' and states that later minimax and moderated-MMD analyses attribute much of the kernel-MMD parametric--nonparametric detection gap to test suboptimality, while the general attribution remains model-dependent. This is supplied dated curation, not independently verified here; the record itself requests independent verification of time-sensitive status and cited claims.
\item The supplied statement status is \texttt{corrected\_verified}. Its verification note reports that the official passage was checked and rewritten as a self-contained problem while preserving its mathematical content. This is a supplied verification claim, not an independent source check performed for this draft.
\end{itemize}
\end{document}
